\documentclass[a4paper,12pt]{article}

\usepackage[top=1cm, bottom=1.8cm, left=1cm, right=1cm]{geometry}
\usepackage[fleqn]{amsmath}
\usepackage{amssymb}
\setlength{\mathindent}{2em}
\usepackage{fontspec}
\usepackage{enumitem}

% Set fonts
\setmainfont{Times New Roman}
\setsansfont{Arial}
\setmonofont{JetBrains Mono}

% Chinese font support
\usepackage{xeCJK}
\setCJKmainfont{Iansui}[
  BoldFont=Iansui,
  AutoFakeBold=3.17
]

\pagestyle{plain}

% Vertical space for student work
\newcommand{\workspace}{\vfill}
\newcommand{\examdate}{2026/03/06}

\begin{document}

\begin{center}
{\Large\bfseries 數學試題}
\end{center}

\begin{flushright}
\examdate
\end{flushright}

\vspace{1em}

\begin{enumerate}[label=\textbf{\arabic*.}, leftmargin=2em, itemsep=0pt]

\item 解方程式：$2x^2 - 5x - 3 = 0$
\workspace

\item 解方程式：$x^2 - 6x + 5 = 0$
\workspace

\item 將下式配方成「完全平方 $+$ 常數」：$x^2 - 8x + 3$
\workspace

\newpage

\item 將下式配方成「完全平方 $+$ 常數」：$3x^2 + 12x - 7$
\workspace

\item 已知點 $P(2,-1)$ 經平移向量 $(-3,5)$ 後得到點 $P'$。求 $P'$ 的座標。
\workspace

\item 將曲線 $(x-4)^2 + (y+1)^2 = 9$ 視為某個以原點為圓心的圓經平移得到。寫出：

\begin{enumerate}[label=(\arabic*)]
\item 原本圓的方程式
\item 平移向量 $(h,k)$
\end{enumerate}
\workspace

\newpage

\item 解聯立方程式：
\[
\begin{cases}
2x + 3y = 7 \\
x - y = 1
\end{cases}
\]
\workspace

\item 解聯立方程式：
\[
\begin{cases}
3x - 2y = 4 \\
5x + y = 13
\end{cases}
\]
\workspace

\end{enumerate}

\end{document}
